% ============================================================
%  GUÍA 2: ARITMÉTICA, ÁLGEBRA Y ECUACIONES
%  Curso Introductorio — 3er Año
%  Autor: Ing. Gabriel Astudillo
%  Sesiones cubiertas: 4, 5, 6
% ============================================================

\documentclass[12pt, a4paper]{article}

% ── Paquetes ─────────────────────────────────────────────────

\usepackage{mathwizards-guia}
\mwguia{Guía 2 — Aritmética, Álgebra y Ecuaciones}{Ing. Gabriel Astudillo $\cdot$ Curso 3er Año}

\begin{document}
% ============================================================

% ── Portada / Título ─────────────────────────────────────────
\mwportada{Guía de Estudio 2}{Aritmética, Álgebra y Ecuaciones}{Números Reales, Exponentes, Notación Científica y Ecuaciones Lineales}

\vspace{4pt}
\begin{cuadroinfo}
\small
\textbf{Presentación:} En esta segunda guía pondremos a punto las herramientas matemáticas fundamentales. Dominar las operaciones con fracciones, entender las leyes de los exponentes y resolver ecuaciones son habilidades que usarás \emph{todo el tiempo} en Física y Química. Piensa en esta guía como el taller donde afilas tus instrumentos antes de la verdadera aventura científica.
\\[4pt]
\textbf{Nivel de Dificultad:}\quad
\facil{} Básico / Directo \quad
\medio{} Intermedio \quad
\dificil{} Desafiante
\end{cuadroinfo}

\vspace{8pt}

% ============================================================
\section{El Conjunto de los Números Reales ($\mathbb{R}$)}
% ============================================================

\subsection{La gran familia de los números}

Los números que usamos a diario forman parte de una gran familia organizada en subconjuntos. Cada subconjunto más grande contiene al anterior:

\begin{cuadroteoria}
\begin{itemize}[leftmargin=*]
  \item \textbf{Números Naturales} ($\mathbb{N}$): Los que usamos para contar: $\{1, 2, 3, 4, \ldots\}$. Algunos autores incluyen el cero.
  \item \textbf{Números Enteros} ($\mathbb{Z}$): Naturales más sus negativos y el cero: $\{\ldots, -3, -2, -1, 0, 1, 2, 3, \ldots\}$.
  \item \textbf{Números Racionales} ($\mathbb{Q}$): Todo número que se puede escribir como fracción $\frac{a}{b}$ con $a, b \in \mathbb{Z}$ y $b \neq 0$. Incluye decimales exactos ($0{,}5 = \frac{1}{2}$) y periódicos ($0{,}333\ldots = \frac{1}{3}$).
  \item \textbf{Números Irracionales} ($\mathbb{I}$): Números que \textbf{no} se pueden expresar como fracción. Su expansión decimal es infinita y no periódica. Ejemplos: $\pi \approx 3{,}14159\ldots$, $\sqrt{2} \approx 1{,}41421\ldots$, el número de oro $\Phi \approx 1{,}61803\ldots$
  \item \textbf{Números Reales} ($\mathbb{R}$): La unión de racionales e irracionales. Cubre toda la recta numérica sin huecos.
\end{itemize}
\end{cuadroteoria}

\begin{figure}[h!]
  \centering
  \begin{tikzpicture}[x=1.5cm, y=1cm]
    % Draw the real line
    \draw[Latex-Latex, thick] (-4.5,0) -- (4.5,0) node[right] {$\mathbb{R}$};
    
    % Draw integer ticks
    \foreach \x in {-4,-3,-2,-1,0,1,2,3,4} {
      \draw[thick] (\x, 0.1) -- (\x, -0.1) node[below, font=\small] {$\x$};
    }
    
    % Draw marked numbers
    % -3
    \filldraw[mwrojonaranja] (-3,0) circle (3pt) node[above=4pt, font=\small\bfseries] {$-3$};
    
    % 1/2
    \filldraw[mwprimario] (0.5,0) circle (3pt) node[above=4pt, font=\small\bfseries] {$\frac{1}{2}$};
    
    % \sqrt{2} ~ 1.414
    \filldraw[mwmagenta] (1.414,0) circle (3pt) node[above=4pt, font=\small\bfseries] {$\sqrt{2}$};
    
    % \pi ~ 3.1416
    \filldraw[mwnaranja] (3.1416,0) circle (3pt) node[above=4pt, font=\small\bfseries] {$\pi$};
    
    % Lines pointing to set labels
    \draw[dashed, thin, gray] (-3,-0.1) -- (-3,-0.7);
    \draw[dashed, thin, gray] (0.5,-0.1) -- (0.5,-0.7);
    \draw[dashed, thin, gray] (1.414,-0.1) -- (1.414,-0.7);
    \draw[dashed, thin, gray] (3.1416,-0.1) -- (3.1416,-0.7);

    % Label sets below
    \node[anchor=north, text=mwprimario, font=\footnotesize] at (0.5,-0.8) {$\frac{1}{2} \in \mathbb{Q}$};
    \node[anchor=north, text=mwrojonaranja, font=\footnotesize] at (-3,-0.8) {$-3 \in \mathbb{Z}$};
    \node[anchor=north, text=mwmagenta, font=\footnotesize] at (1.414,-0.8) {$\sqrt{2} \in \mathbb{I}$};
    \node[anchor=north, text=mwnaranja, font=\footnotesize] at (3.1416,-0.8) {$\pi \in \mathbb{I}$};
    
    % General subset relation label
    \node[draw=mwprimario, fill=mwprimario!12, rounded corners, inner sep=4pt, font=\footnotesize\bfseries, text=mwprimario] at (0,1.5) {$\mathbb{N} \subset \mathbb{Z} \subset \mathbb{Q} \subset \mathbb{R}$};
  \end{tikzpicture}
  \caption{La recta numérica y los conjuntos numéricos.}
  \label{fig:recta}
\end{figure}

\begin{cuadroadvertencia}
\textbf{Dato curioso — El número de oro ($\Phi$):} Se define como $\Phi = \frac{1 + \sqrt{5}}{2} \approx 1{,}618$. Aparece en la naturaleza con sorprendente frecuencia: en la disposición de los pétalos de las flores, en la espiral de las conchas marinas e incluso en las proporciones del cuerpo humano. Los griegos lo llamaban «la proporción divina».
\end{cuadroadvertencia}

\subsection{Repaso intensivo de fracciones}

Las fracciones son la base de todo el álgebra. Si dominas las cuatro operaciones con fracciones, el despeje de fórmulas en Física será pan comido. Repasemos:

\begin{cuadrosolucion}
\textbf{Operaciones fundamentales con fracciones:}
\begin{enumerate}[leftmargin=*, label=\arabic*.]
  \item \textbf{Suma y resta} (mismo denominador): $\dfrac{a}{c} \pm \dfrac{b}{c} = \dfrac{a \pm b}{c}$
  
  \item \textbf{Suma y resta} (distinto denominador): Se busca el mínimo común múltiplo (mcm) de los denominadores.
  $$\frac{a}{b} + \frac{c}{d} = \frac{a \cdot d + c \cdot b}{b \cdot d}$$
  
  \item \textbf{Multiplicación}: Se multiplican numeradores entre sí y denominadores entre sí.
  $$\frac{a}{b} \cdot \frac{c}{d} = \frac{a \cdot c}{b \cdot d}$$
  
  \item \textbf{División}: Se multiplica por el recíproco de la segunda fracción.
  $$\frac{a}{b} \div \frac{c}{d} = \frac{a}{b} \cdot \frac{d}{c} = \frac{a \cdot d}{b \cdot c}$$
\end{enumerate}
\end{cuadrosolucion}

\noindent\textbf{Ejemplo resuelto:} Simplifica $\dfrac{3}{4} + \dfrac{2}{5} - \dfrac{1}{10}$.
\begin{itemize}
  \item mcm$(4, 5, 10) = 20$.
  \item $\dfrac{3}{4} = \dfrac{15}{20}$, \quad $\dfrac{2}{5} = \dfrac{8}{20}$, \quad $\dfrac{1}{10} = \dfrac{2}{20}$.
  \item Resultado: $\dfrac{15 + 8 - 2}{20} = \dfrac{21}{20}$.
\end{itemize}

\newpage

% ============================================================
\subsection{Ejercicios: Números Reales y Fracciones}
% ============================================================

\begin{enumerate}[leftmargin=*]
  \item \facil{} Clasifica cada número en el conjunto más pequeño al que pertenece ($\mathbb{N}$, $\mathbb{Z}$, $\mathbb{Q}$ o $\mathbb{I}$):
  \begin{multicols}{3}
  \begin{enumerate}[label=\alph*)]
    \item $-7$
    \item $\frac{5}{3}$
    \item $\sqrt{9}$
    \item $0{,}1\overline{6}$
    \item $\pi$
    \item $0$
  \end{enumerate}
  \end{multicols}

  \item \facil{} Calcula: $\dfrac{2}{3} + \dfrac{5}{6}$

  \item \facil{} Calcula: $\dfrac{7}{8} - \dfrac{3}{4}$

  \item \facil{} Calcula: $\dfrac{4}{9} \times \dfrac{3}{2}$

  \item \facil{} Calcula: $\dfrac{5}{6} \div \dfrac{10}{3}$

  \item \medio{} Simplifica: $\dfrac{3}{4} + \dfrac{2}{5} - \dfrac{1}{2}$

  \item \medio{} Resuelve: $\dfrac{2}{3} \times \left(\dfrac{5}{4} - \dfrac{1}{2}\right)$

  \item \medio{} Un tanque está lleno hasta $\frac{3}{5}$ de su capacidad. Si se vacía $\frac{1}{4}$ del tanque completo, ¿qué fracción del tanque queda con agua?

  \item \medio{} Ordena de menor a mayor: $\dfrac{3}{5}$, $\dfrac{7}{10}$, $\dfrac{2}{3}$, $\dfrac{5}{8}$.

  \item \dificil{} Simplifica: $\dfrac{\frac{2}{3} + \frac{1}{4}}{\frac{5}{6} - \frac{1}{3}}$

  \item \dificil{} Si $\frac{2}{7}$ de un número es igual a 18, ¿cuál es el número? ¿Y cuál es $\frac{3}{7}$ de ese mismo número?

  \item \dificil{} Demuestra que $0{,}999\ldots = 1$ usando álgebra. (Pista: sea $x = 0{,}999\ldots$, multiplica por 10 y resta.)
\end{enumerate}

\newpage

% ============================================================
\section{Exponentes y Notación Científica}
% ============================================================

\subsection{Potenciación y sus propiedades}

La potenciación es una forma abreviada de escribir multiplicaciones repetidas. Si multiplicamos la base $a$ por sí misma $n$ veces:
$$a^n = \underbrace{a \times a \times a \times \cdots \times a}_{n \text{ veces}}$$

\begin{cuadroteoria}
\textbf{Propiedades fundamentales de los exponentes} (con $a, b \neq 0$):
\begin{enumerate}[leftmargin=*, label=\arabic*.]
  \item $a^m \cdot a^n = a^{m+n}$ \hfill (Producto de potencias de igual base)
  \item $\dfrac{a^m}{a^n} = a^{m-n}$ \hfill (Cociente de potencias de igual base)
  \item $(a^m)^n = a^{m \cdot n}$ \hfill (Potencia de una potencia)
  \item $(a \cdot b)^n = a^n \cdot b^n$ \hfill (Potencia de un producto)
  \item $\left(\dfrac{a}{b}\right)^n = \dfrac{a^n}{b^n}$ \hfill (Potencia de un cociente)
  \item $a^0 = 1$ \hfill (Todo número no nulo elevado a cero es 1)
  \item $a^{-n} = \dfrac{1}{a^n}$ \hfill (Exponente negativo $=$ recíproco)
\end{enumerate}
\end{cuadroteoria}

\noindent\textbf{Ejemplo:} Simplifica $\dfrac{2^5 \cdot 2^3}{2^4}$.
\begin{itemize}
  \item Numerador: $2^5 \cdot 2^3 = 2^{5+3} = 2^8$ (propiedad 1).
  \item Cociente: $\dfrac{2^8}{2^4} = 2^{8-4} = 2^4 = 16$ (propiedad 2).
\end{itemize}

\subsection{Notación Científica}

En ciencias, trabajamos con números extremadamente grandes o extremadamente pequeños. La \textbf{notación científica} nos permite expresarlos de forma compacta:

\begin{cuadrosolucion}
$$N = a \times 10^n$$
Donde:
\begin{itemize}[leftmargin=*]
  \item $a$ es un número tal que $1 \leq |a| < 10$ (tiene \textbf{un solo dígito} antes del punto decimal).
  \item $n$ es un número entero que indica cuántas posiciones se «movió» el punto decimal.
  \item Si $n > 0$: el número original es \textbf{grande}, al contrario, si $n < 0$: el número original es \textbf{pequeño}.
\end{itemize}
\end{cuadrosolucion}

\noindent\textbf{Ejemplos:}
\begin{itemize}
  \item Distancia de la Tierra al Sol: $149\,600\,000\,\text{km} = 1{,}496 \times 10^8\,\text{km}$
  \item Diámetro de un átomo de hidrógeno: $0{,}000\,000\,000\,106\,\text{m} = 1{,}06 \times 10^{-10}\,\text{m}$
  \item Velocidad de la luz: $300\,000\,000\,\text{m/s} \approx 3 \times 10^8\,\text{m/s}$
\end{itemize}

\begin{cuadroadvertencia}
\textbf{Truco para recordar:} Si mueves el punto decimal hacia la \emph{izquierda}, el exponente es \textbf{positivo}. Si lo mueves hacia la \emph{derecha}, es \textbf{negativo}. Cuenta la cantidad de saltos y ese es el valor de $n$.
\end{cuadroadvertencia}

%\newpage

% ============================================================
\subsection{Ejercicios: Exponentes y Notación Científica}
% ============================================================

\begin{enumerate}[leftmargin=*, resume]
  \item \facil{} Simplifica usando las propiedades de los exponentes: $3^4 \cdot 3^2$

  \item \facil{} Simplifica: $\dfrac{5^7}{5^3}$

  \item \facil{} Calcula: $(2^3)^2$

  \item \facil{} Expresa en notación científica: $47\,000\,000$

  \item \medio{} Simplifica: $\dfrac{4^3 \cdot 4^{-1}}{4^2}$

  \item \medio{} Expresa en notación científica: $0{,}000\,032$

  \item \medio{} Calcula y expresa en notación científica: $(3 \times 10^4) \times (2 \times 10^5)$

  \item \medio{} Calcula y expresa en notación científica: $\dfrac{8 \times 10^9}{4 \times 10^3}$

  \item \medio{} Simplifica completamente: $\dfrac{(2^3)^2 \cdot 3^4}{2^4 \cdot 3^2}$

  \item \dificil{} La masa de la Tierra es aproximadamente $5{,}97 \times 10^{24}$\,kg y la masa de la Luna es $7{,}35 \times 10^{22}$\,kg. ¿Cuántas veces más masiva es la Tierra que la Luna?

  \item \dificil{} Simplifica: $\left(\dfrac{2^{-3} \cdot 4^2}{8}\right)^{-1}$. (Pista: escribe todo en base 2.)

  \item \dificil{} Una bacteria mide $2 \times 10^{-6}$\,m. Si se colocaran $5 \times 10^3$ bacterias en fila, ¿qué longitud total tendrían en metros y en milímetros?
\end{enumerate}

\newpage

% ============================================================
\section{Ecuaciones Lineales y Sistemas de Ecuaciones}
% ============================================================

\subsection{La igualdad como balanza}

Una \textbf{ecuación} es una igualdad que contiene una o más incógnitas. Resolverla significa encontrar el valor de la incógnita que hace que ambos lados de la igualdad tengan el mismo valor.

Piensa en una ecuación como una balanza en equilibrio: lo que haces de un lado, debes hacerlo también del otro para mantener el equilibrio.

\begin{figure}[h!]
  \centering
  \begin{tikzpicture}[x=1cm, y=1cm]
    % Central Support
    \draw[thick, fill=gray!30] (0,-1) -- (0.3,-2.5) -- (-0.3,-2.5) -- cycle; % Base stand
    \draw[ultra thick, gray!80] (0,-1) -- (0,-0.3); % Stand post
    \filldraw[gray] (0,-0.3) circle (3pt); % Pivot point
    
    % Main Beam (balanced)
    \draw[ultra thick, mwprimario] (-3,-0.3) -- (3,-0.3);
    
    % Hanger strings (left side)
    \draw[thick] (-3,-0.3) -- (-3.5,-1.3);
    \draw[thick] (-3,-0.3) -- (-2.5,-1.3);
    % Plate (left side)
    \draw[very thick, fill=mwprimario!12, draw=mwprimario] (-3.7,-1.3) -- (-2.3,-1.3) -- (-2.5,-1.4) -- (-3.5,-1.4) -- cycle;
    \node[text=mwprimario, font=\bfseries] at (-3,-1.8) {$3x + 2$};
    
    % Hanger strings (right side)
    \draw[thick] (3,-0.3) -- (2.5,-1.3);
    \draw[thick] (3,-0.3) -- (3.5,-1.3);
    % Plate (right side)
    \draw[very thick, fill=mwprimario!12, draw=mwprimario] (2.3,-1.3) -- (3.7,-1.3) -- (3.5,-1.4) -- (2.5,-1.4) -- cycle;
    \node[text=mwprimario, font=\bfseries] at (3,-1.8) {$14$};
    
    % Equal sign in the middle
    \node[text=mwmagenta, font=\Huge\bfseries] at (0,-1.8) {$=$};
  \end{tikzpicture}
  \caption{La ecuación como balanza: lo que haces de un lado, lo haces del otro.}
  \label{fig:balanza}
\end{figure}

\begin{cuadroteoria}
\textbf{Principios fundamentales para resolver ecuaciones:}
\begin{enumerate}[leftmargin=*, label=\arabic*.]
  \item \textbf{Principio de adición:} Si $a = b$, entonces $a + c = b + c$ y $a - c = b - c$.
  
  \textit{Puedes sumar o restar la misma cantidad a ambos lados.}
  
  \item \textbf{Principio de multiplicación:} Si $a = b$, entonces $a \cdot c = b \cdot c$ y $\frac{a}{c} = \frac{b}{c}$ (con $c \neq 0$).
  
  \textit{Puedes multiplicar o dividir ambos lados por la misma cantidad (distinta de cero).}
\end{enumerate}
\end{cuadroteoria}

\subsection{Resolución de ecuaciones lineales}

Una ecuación lineal tiene la forma $ax + b = c$, donde la incógnita $x$ aparece con exponente 1.

\noindent\textbf{Ejemplo resuelto:} Resuelve $5x - 3 = 2x + 9$.
\begin{enumerate}
  \item Agrupar términos con $x$ a un lado: $5x - 2x = 9 + 3$
  \item Simplificar: $3x = 12$
  \item Dividir por el coeficiente: $x = \frac{12}{3} = 4$
  \item \textbf{Verificación:} $5(4) - 3 = 17$ y $2(4) + 9 = 17$ $\checkmark$
\end{enumerate}

\newpage

\subsection{Sistemas de ecuaciones lineales}

Cuando tenemos \textbf{dos incógnitas}, necesitamos \textbf{dos ecuaciones} para encontrar una solución única. Los dos métodos principales son:

\begin{cuadrosolucion}
\textbf{Método de sustitución:}
\begin{enumerate}[leftmargin=*]
  \item Despeja una variable de una ecuación.
  \item Sustituye esa expresión en la otra ecuación.
  \item Resuelve la ecuación resultante (tiene una sola incógnita).
  \item Sustituye el valor hallado para encontrar la segunda variable.
\end{enumerate}
\end{cuadrosolucion}

\begin{cuadroextra}
\textbf{Método de reducción (eliminación):}
\begin{enumerate}[leftmargin=*]
  \item Multiplica una o ambas ecuaciones por constantes para que una variable tenga coeficientes opuestos.
  \item Suma las ecuaciones para eliminar esa variable.
  \item Resuelve la ecuación resultante.
  \item Sustituye para hallar la otra variable.
\end{enumerate}
\end{cuadroextra}

\noindent\textbf{Ejemplo resuelto (Sustitución):}
$$\begin{cases} x + y = 10 \\ 2x - y = 5 \end{cases}$$

\begin{itemize}
  \item De la primera: $y = 10 - x$.
  \item Sustituimos en la segunda: $2x - (10 - x) = 5 \Rightarrow 3x - 10 = 5 \Rightarrow 3x = 15 \Rightarrow x = 5$.
  \item Entonces $y = 10 - 5 = 5$.
  \item \textbf{Verificación:} $5 + 5 = 10$ $\checkmark$ y $2(5) - 5 = 5$ $\checkmark$.
\end{itemize}

\newpage

% ============================================================
\subsection{Ejercicios: Ecuaciones Lineales y Sistemas}
% ============================================================

\begin{enumerate}[leftmargin=*, resume]
  \item \facil{} Resuelve: $3x + 7 = 22$

  \item \facil{} Resuelve: $4x - 5 = 2x + 9$

  \item \facil{} Resuelve: $\dfrac{x}{3} = 8$

  \item \facil{} Resuelve: $2(x + 3) = 18$

  \item \medio{} Resuelve: $\dfrac{2x - 1}{3} = \dfrac{x + 4}{2}$

  \item \medio{} Resuelve el sistema por sustitución:
  $\begin{cases} x + y = 15 \\ x - y = 3 \end{cases}$

  \item \medio{} Resuelve el sistema por reducción:
  $\begin{cases} 2x + 3y = 12 \\ 4x - 3y = 6 \end{cases}$

  \item \medio{} Resuelve: $5(x - 2) + 3 = 2(x + 4) - 1$

  \item \dificil{} Un número y su tercera parte suman 20. Plantea la ecuación y resuélvela.

  \item \dificil{} Resuelve el sistema:
  $\begin{cases} 3x + 2y = 16 \\ 5x - 4y = -2 \end{cases}$

  \item \dificil{} En una tienda, 3 cuadernos y 2 lápices cuestan 28\,Bs, mientras que 5 cuadernos y 3 lápices cuestan 45\,Bs. Plantea el sistema de ecuaciones y halla el precio de cada artículo.

  \item \dificil{} Resuelve: $\dfrac{x + 2}{4} - \dfrac{x - 1}{3} = \dfrac{5}{6}$
\end{enumerate}

\newpage

% ============================================================
\section{Clave de Respuestas}
% ============================================================

\subsection*{Números Reales y Fracciones (Ejercicios 1--12)}

\begin{enumerate}[leftmargin=*, label=\textbf{\arabic*.}]
  \item a) $-7 \in \mathbb{Z}$. \quad b) $\frac{5}{3} \in \mathbb{Q}$. \quad c) $\sqrt{9} = 3 \in \mathbb{N}$. \quad d) $0{,}1\overline{6} = \frac{1}{6} \in \mathbb{Q}$. \quad e) $\pi \in \mathbb{I}$. \quad f) $0 \in \mathbb{Z}$ (o $\mathbb{N}$ según convención).

  \item $\frac{2}{3} + \frac{5}{6} = \frac{4}{6} + \frac{5}{6} = \frac{9}{6} = \frac{3}{2}$

  \item $\frac{7}{8} - \frac{3}{4} = \frac{7}{8} - \frac{6}{8} = \frac{1}{8}$

  \item $\frac{4}{9} \times \frac{3}{2} = \frac{12}{18} = \frac{2}{3}$

  \item $\frac{5}{6} \div \frac{10}{3} = \frac{5}{6} \times \frac{3}{10} = \frac{15}{60} = \frac{1}{4}$

  \item $\frac{3}{4} + \frac{2}{5} - \frac{1}{2} = \frac{15}{20} + \frac{8}{20} - \frac{10}{20} = \frac{13}{20}$

  \item $\frac{2}{3} \times \left(\frac{5}{4} - \frac{1}{2}\right) = \frac{2}{3} \times \left(\frac{5}{4} - \frac{2}{4}\right) = \frac{2}{3} \times \frac{3}{4} = \frac{6}{12} = \frac{1}{2}$

  \item $\frac{3}{5} - \frac{1}{4} = \frac{12}{20} - \frac{5}{20} = \frac{7}{20}$ del tanque.

  \item Convertimos a denominador común (40): $\frac{3}{5} = \frac{24}{40}$, $\frac{5}{8} = \frac{25}{40}$, $\frac{2}{3} = \frac{26{,}\overline{6}}{40}$... Mejor con 30: $\frac{3}{5} = \frac{18}{30}$, $\frac{7}{10} = \frac{21}{30}$, $\frac{2}{3} = \frac{20}{30}$, $\frac{5}{8} = \frac{18{,}75}{30}$. Orden: $\frac{3}{5} < \frac{5}{8} < \frac{2}{3} < \frac{7}{10}$.

  \item $\frac{\frac{2}{3} + \frac{1}{4}}{\frac{5}{6} - \frac{1}{3}} = \frac{\frac{8+3}{12}}{\frac{5-2}{6}} = \frac{\frac{11}{12}}{\frac{3}{6}} = \frac{11}{12} \times \frac{6}{3} = \frac{11}{12} \times 2 = \frac{11}{6}$

  \item $\frac{2}{7}x = 18 \Rightarrow x = 18 \times \frac{7}{2} = 63$. Entonces $\frac{3}{7} \times 63 = 27$.

  \item Sea $x = 0{,}999\ldots$ Entonces $10x = 9{,}999\ldots$ Restando: $10x - x = 9{,}999\ldots - 0{,}999\ldots \Rightarrow 9x = 9 \Rightarrow x = 1$.
\end{enumerate}

\subsection*{Exponentes y Notación Científica (Ejercicios 13--24)}

\begin{enumerate}[leftmargin=*, label=\textbf{\arabic*.}]
  \setcounter{enumi}{12}
  \item $3^4 \cdot 3^2 = 3^{4+2} = 3^6 = 729$

  \item $\frac{5^7}{5^3} = 5^{7-3} = 5^4 = 625$

  \item $(2^3)^2 = 2^{3 \times 2} = 2^6 = 64$

  \item $47\,000\,000 = 4{,}7 \times 10^7$

  \item $\frac{4^3 \cdot 4^{-1}}{4^2} = \frac{4^{3-1}}{4^2} = \frac{4^2}{4^2} = 4^0 = 1$

  \item $0{,}000\,032 = 3{,}2 \times 10^{-5}$

  \item $(3 \times 10^4)(2 \times 10^5) = 6 \times 10^{4+5} = 6 \times 10^9$

  \item $\frac{8 \times 10^9}{4 \times 10^3} = 2 \times 10^{9-3} = 2 \times 10^6$

  \item $\frac{(2^3)^2 \cdot 3^4}{2^4 \cdot 3^2} = \frac{2^6 \cdot 3^4}{2^4 \cdot 3^2} = 2^{6-4} \cdot 3^{4-2} = 2^2 \cdot 3^2 = 4 \cdot 9 = 36$

  \item $\frac{5{,}97 \times 10^{24}}{7{,}35 \times 10^{22}} = \frac{5{,}97}{7{,}35} \times 10^{24-22} = 0{,}8122 \times 10^2 \approx 81{,}2$ veces.

  \item En base 2: $2^{-3} = 2^{-3}$, $4^2 = (2^2)^2 = 2^4$, $8 = 2^3$. Interior: $\frac{2^{-3} \cdot 2^4}{2^3} = \frac{2^1}{2^3} = 2^{-2}$. Con exponente $-1$: $(2^{-2})^{-1} = 2^2 = 4$.

  \item $2 \times 10^{-6} \times 5 \times 10^3 = 10 \times 10^{-3} = 1 \times 10^{-2}\,\text{m} = 0{,}01\,\text{m} = 10\,\text{mm}$.
\end{enumerate}

\subsection*{Ecuaciones Lineales y Sistemas (Ejercicios 25--36)}

\begin{enumerate}[leftmargin=*, label=\textbf{\arabic*.}]
  \setcounter{enumi}{24}
  \item $3x + 7 = 22 \Rightarrow 3x = 15 \Rightarrow x = 5$.

  \item $4x - 5 = 2x + 9 \Rightarrow 2x = 14 \Rightarrow x = 7$.

  \item $\frac{x}{3} = 8 \Rightarrow x = 24$.

  \item $2(x+3) = 18 \Rightarrow 2x + 6 = 18 \Rightarrow 2x = 12 \Rightarrow x = 6$.

  \item $\frac{2x-1}{3} = \frac{x+4}{2}$. Multiplicamos por 6: $2(2x-1) = 3(x+4) \Rightarrow 4x - 2 = 3x + 12 \Rightarrow x = 14$.

  \item Sumando: $2x = 18 \Rightarrow x = 9$, $y = 15 - 9 = 6$.

  \item Sumando: $6x = 18 \Rightarrow x = 3$. En la primera: $6 + 3y = 12 \Rightarrow y = 2$.

  \item $5x - 10 + 3 = 2x + 8 - 1 \Rightarrow 5x - 7 = 2x + 7 \Rightarrow 3x = 14 \Rightarrow x = \frac{14}{3}$.

  \item $x + \frac{x}{3} = 20 \Rightarrow \frac{4x}{3} = 20 \Rightarrow x = 15$.

  \item Multiplicamos la primera por 2: $6x + 4y = 32$. Sumamos con la segunda: $11x = 30 \Rightarrow x = \frac{30}{11}$. Sustituyendo: $y = \frac{16 - 3 \cdot \frac{30}{11}}{2} = \frac{16 - \frac{90}{11}}{2} = \frac{\frac{176 - 90}{11}}{2} = \frac{86}{22} = \frac{43}{11}$.

  \item Sea $c$ el precio del cuaderno y $l$ el del lápiz. Sistema: $3c + 2l = 28$ y $5c + 3l = 45$. Multiplicamos la primera por 3 y la segunda por $-2$: $9c + 6l = 84$, $-10c - 6l = -90$. Sumando: $-c = -6 \Rightarrow c = 6$. Sustituyendo: $18 + 2l = 28 \Rightarrow l = 5$. Cuaderno: 6\,Bs, lápiz: 5\,Bs.

  \item Multiplicamos por 12: $3(x+2) - 4(x-1) = 10 \Rightarrow 3x + 6 - 4x + 4 = 10 \Rightarrow -x + 10 = 10 \Rightarrow x = 0$.
\end{enumerate}

\end{document}
